\section{Introduction}

The cost of an action measures how many edges, on average, are needed
to connect its orbits. More precisely, let a countable group $\Gamma$
act by measure-preserving Borel automorphisms on a standard probability
space $(X,\mu)$. We use right actions, so $(xg)h=x(gh)$.
Its orbit relation is
\[
 \mathcal R=\{(x,xg):x\in X,\ g\in\Gamma\}.
\]
A \emph{graphing} $\mathcal E$ of $\mathcal R$ is a countable family of
measure-preserving Borel bijections between Borel subsets of $X$,
with graphs contained in $\mathcal R$. It generates $\mathcal R$ if
finite paths in these maps and their inverses connect all related
points on a conull set. Its cost $C(\mathcal E)$ is the sum of its domain measures;
$\operatorname{Cost}(\mathcal R)$ is the infimum of these sums over
generating graphings. An action is \emph{essentially free} if each
nonidentity group element fixes a null set. A group has \emph{fixed
price} if all its essentially free probability-measure-preserving
(p.m.p.) actions have the same cost.

Levitt introduced cost \cite{levitt-1995}. Gaboriau developed its
connections with group theory and the fixed-price problem
\cite{gaboriau-2000}. His work shows, in particular, that free groups
have fixed price equal to their rank, and establishes fixed price one
for many groups with commuting or amenable structure. Recent results
include the fixed-price-one theorem for products of two infinite
countable groups \cite{khezeli-2026-v4}, metric and infinite-measure
action criteria \cite{bevilacqua-bowen-2025-v1}, and a criterion based
on small product neighborhoods of finite subsets
\cite{slutsky-2026-v1}. We give a group for which fixed price fails.

\subsection{The group and the two actions}

Let
\[
 A=F(a,b_1,\ldots,b_{99}),\qquad
 w=ab_1ab_2\cdots ab_{99}a,
 \qquad J=\langle b_1,\ldots,b_{99},w\rangle\le A.
\]
Let $\langle t\rangle$ be infinite cyclic and set
\begin{equation}\label{intro:group}
 \Gamma=A*_J(J\times\langle t\rangle).
\end{equation}
The homomorphism $\chi:\Gamma\to\mathbb Z$ given by
$\chi(a)=1$ and $\chi(b_i)=\chi(t)=0$ is well defined; note that
$\chi(w)=100$. On the Bernoulli
probability space $X_0=[0,1]^\Gamma$ with product Lebesgue measure, let
\[
 (xg)(h)=x(gh).
\]
Let $X\subseteq X_0$ be the invariant conull set of points with trivial
stabilizer, equipped with the restricted probability measure. For an
integer $M>100$, give $Y_M=X\times\mathbb Z/M\mathbb Z$
product measure with uniform height coordinate, and define
\[
 (x,r)g=(xg,r+\chi(g)\bmod M).
\]
Both actions are free; restriction from $X_0$ to $X$ preserves cost.
Their elementary group-theoretic
and measure-theoretic properties are proved in Section~\ref{actions:section}.

\begin{theorem}\label{intro:main}
Put $K=e^{96}96^{99}/95^{95}$, choose
\[
 0<\alpha<1/200,\qquad K\alpha^3<1/2,
 \qquad \eta=\alpha/100.
\]
For the group and actions just defined,
\[
 \operatorname{Cost}(\mathcal R_{\Gamma\curvearrowright X})
      \ge1+\eta,
 \qquad
 \operatorname{Cost}(\mathcal R_{\Gamma\curvearrowright Y_M})
      \le1+\frac{99}{M}\quad(M>100).
\]
In particular, choosing an integer $M>\max\{100,99/\eta\}$ gives
two essentially free p.m.p. actions of the same finitely generated
group with strictly different costs. Thus the general fixed-price
assertion is false.
\end{theorem}

\subsection{How the proof works}

The upper bound uses the finite height coordinate. All $t$-edges,
together with the generators of $J$ on a small set meeting every
$t$-orbit, generate the $J\times\langle t\rangle$ relation at cost
arbitrarily close to one. It remains to add $a$-edges. The relation
$w=ab_1\cdots ab_{99}a$ propagates $a$-edges from any $99$
consecutive heights to the next height, so an initial cost $99/M$
is enough.

For the lower bound, write $\mathcal R_A$ and $\mathcal R_J$ for
the restricted orbit relations on $X$, and define the relative
quantity
\[
 \delta_X=\inf\{C(\mathcal E):
       \mathcal R_J\vee\mathcal E=\mathcal R_A,
       \ \mathcal E\subseteq\mathcal R_A\}.
\]
Here $C(\mathcal E)$ is the graphing cost, and the join denotes the
relation generated by the supplied $\mathcal R_J$-steps and the
edges of $\mathcal E$. Three steps give the desired bound:
\begin{enumerate}
\item \emph{Deploy and compress graphings.}
Subdivide a nearly optimal $\Gamma$-graphing into factor edges on a
larger finite-measure space. Repeated compression makes more of the
$J$-relation available and recovers its exact saved cost. Amalgam
normal form then gives
$\operatorname{Cost}(\mathcal R_\Gamma)\ge1+\delta_X$.
The construction follows Gaboriau's deployment method
\cite[\S\,IV.28--IV.32 and Proposition~IV.35]{gaboriau-2000};
Sections~\ref{compression:section}--\ref{deployment:section}
prove the required relative estimate directly.
\item \emph{Pass to finite permutation models.}
Put $u_0=1$ and $u_i=u_{i-1}ab_i$ for $1\le i\le99$. A finite set $V$ with a right
$A$-action by permutations defines a group whose symbols record
$a$-steps, while all $J$-steps are assigned the identity:
\[
 D_V=\langle x_v\ (v\in V)\mid
              x_vx_{vu_1}\cdots x_{vu_{99}}=1\ (v\in V)\rangle.
\]
Write $\operatorname{rk}(D_V)$ for its minimum number of generators.
Local approximation of Bernoulli graphings gives
$\delta_X\ge\limsup\operatorname{rk}(D_V)/|V|$ along any
asymptotically free sequence of such models, meaning that each fixed
nonidentity element of $A$ fixes a fraction of points tending to zero.
The displayed row makes $w$ trivial in this encoding. In
Section~\ref{models:section}, random permutations supply models with
strong expansion and few overlapping relators.
\item \emph{Bound the rank by shortening a labelled graph.}
The expansion property forces
$\operatorname{rk}(D_V)>\eta|V|$ for large models. A minimal
labelled graph representing a generating set would otherwise contain
a long relator segment that could be replaced by a shorter path.
This uses the Arzhantseva--Ol'shanskii method
\cite{arzhantseva-olshanskii-1996,kapovich-schupp-2002}.
The additional ingredients are a saturation procedure for exceptional
letters and a planar word lemma over an arbitrary coefficient group,
proved in Sections~\ref{planar:section} and~\ref{rank:section}.
\end{enumerate}

The proof isolates three estimates beyond this example: compression
relative to a supplied equivalence relation, a transfer from Bernoulli
graphings to ranks of finitely presented groups, and a deterministic
rank bound from expansion with sparse overlaps. Their statements
specify their own inputs and conclusions. The coefficient group is the actual subgroup generated by the letters
absorbed during saturation. It need not be finitely presented: all
relations true in that subgroup are retained, so its injection into the
presented group is part of the construction. The planar lemma keeps
track of these coefficient values while counting only the remaining
letters.
